% ****** Start of file apssamp.tex ******
%
%   This file is part of the APS files in the REVTeX 4.2 distribution.
%   Version 4.2a of REVTeX, December 2014
%
%   Copyright (c) 2014 The American Physical Society.
%
%   See the REVTeX 4 README file for restrictions and more information.
%
% TeX'ing this file requires that you have AMS-LaTeX 2.0 installed
% as well as the rest of the prerequisites for REVTeX 4.2
%
% See the REVTeX 4 README file
% It also requires running BibTeX. The commands are as follows:
%
%  1)  latex apssamp.tex
%  2)  bibtex apssamp
%  3)  latex apssamp.tex
%  4)  latex apssamp.tex
%
\documentclass[%
 reprint,
%superscriptaddress,
%groupedaddress,
%unsortedaddress,
%runinaddress,
%frontmatterverbose, 
%preprint,
%preprintnumbers,
%nofootinbib,
%nobibnotes,
%bibnotes,
 amsmath,amssymb,
 aps,
%pra,
%prb,
%rmp,
%prstab,
%prstper,
%floatfix,
]{revtex4-2}

\usepackage{graphicx}% Include figure files
\usepackage{dcolumn}%
\usepackage{physics}

%Align table columns on decimal point
\usepackage{bm}% bold math
%\usepackage{hyperref}% add hypertext capabilities
%\usepackage[mathlines]{lineno}% Enable numbering of text and display math
%\linenumbers\relax % Commence numbering lines

\usepackage{amsmath}

%\usepackage[showframe,%Uncomment any one of the following lines to test 
%%scale=0.7, marginratio={1:1, 2:3}, ignoreall,% default settings
%%text={7in,10in},centering,
%%margin=1.5in,
%%total={6.5in,8.75in}, top=1.2in, left=0.9in, includefoot,
%%height=10in,a5paper,hmargin={3cm,0.8in},
%]{geometry}

\begin{document}

\preprint{APS/123-QED}

\title{Superior Performance of Phase Binarized Oscillators (PBOs) \\ Compared to Quantum Approximation Optimization Algorithm (QAOA) \\ for Ising Computing (Max-Cut Problem)}% Force line breaks with \\


\author{Sanyam Singhal}
 \affiliation{Department of Physics, Indian Institute of Technology Bombay, Maharashtra 400076, India}%Lines break automatically or can be forced with \\
\author{Debanjan Bhowmik}%
 \email{debanjan@ee.iitb.ac.in}
\affiliation{Department of Electrical Engineering, Indian Institute of Technology Bombay, Maharashtra 400076, India
}%



\date{\today}% It is always \today, today,
             %  but any date may be explicitly specified

\begin{abstract}
It has been shown both theoretically and experimentally that most coupled oscillators undergo synchronization and phase binarization at room temperature under sub-harmonic injection locking (SHIL), irrespective of the exact physics of their auto-oscillation or coupling. These phase-binarized oscillators (PBOs) can be used for Ising computing, which is about heuristically solving NP-Hard combinatorial optimization problems very fast. The quantum approximate optimization algorithm (QAOA) has emerged as an alternative noisy intermediate scale quantum (NISQ) era algorithm for Ising computing and is very popular currently since it is gate based and needs low circuit depth for implementation. In this paper, we compare the performance of PBOs with that of QAOA over a wide range of graph instances of different levels of difficulty, while restricting ourselves to the NP-Hard Max-Cut problem. We show that for the difficult graph instances (unweighted random cubic, unweighted Erd{\"o}s R{\'e}nyi, and weighted complete graphs with relatively high number of nodes: 18-20), the success probability (probability to find the correct Max-Cut solution) of PBOs is 4-5 orders of magnitude higher than that of QAOA. In fact, for Erd{\"o}s R{\'e}nyi and complete graphs, success probability of QAOA drops exponentially with increase in the number of nodes of the graph (graph size), while success probability of PBO only drops linearly with the number of nodes. Since PBOs operate at room temperature while the quantum circuit in QAOA doesn't (it operates in milli Kelvins), our finding here on their success probability numbers makes PBOs a more attractive hardware platform for Ising computing compared to quantum approaches like QAOA. Since we use the very general and physics-agnostic Kuramoto model to simulate PBOs here, our result is applicable to a wide range of oscillators both based on conventional transistors and emerging nanoscale devices. Hence, we also compare the time to solution for PBOs based on these different device technologies in this paper.
\end{abstract}

%\keywords{Suggested keywords}%Use showkeys class option if keyword
                              %display desired
\maketitle

%\tableofcontents

\section{Introduction}

Coupled oscillators are known to show very interesting phase dynamics, the nature of which is often independent of the microscopic physics of the auto-oscillation or the coupling. As a result, generalized mathematical models like the Kuramoto model \citep{Kuramoto} and its variations (Kuramoto-Sakaguchi model \citep{Circuit5,Circuit6}, generalized Adler's equation, etc. \citep{Jaijeet1,Jaijeet2}) have been used to describe the dynamics. It has been shown analytically, as well as through numeric simulations and experiments, that the phase dynamics of various kinds of oscillators, both based on conventional electronic circuits \citep{Circuit5,Circuit6} as well as emerging nano devices like metal-to-insulator transition devices \citep{MIT1,MIT2} and spintronic devices \citep{Spin_Liu,Spin_DB1,Spin_DB2,Spin_Ashwin,Spin_Hyunsoo,Spin_Jullie2,Spin_Finocchio1,Spin_Finocchio2}, can be captured by these models. Subsequently, it has also been shown that these oscillators, even at room temperature, undergo phase binarization upon sub-harmonic injection locking (SHIL) \citep{Jaijeet1,Jaijeet2,Spin_Akerman1,Spin_Akerman3}. These phase-binarized oscillators (PBOs) have, in turn, been shown (both through experiments and simulations) to carry out Ising computing, which is about heuristically solving NP-Hard combinatorial optimization problems much faster than conventional computers can \citep{IsingComputingReview,Jaijeet2,Circuit1,Circuit2,Circuit3,Circuit4,Spin_Finocchio2,Spin_Akerman3}. Solving NP-hard combinatorial optimization problems is not only of interest to the theoretical computer science community but is also useful for various practical applications: delivery scheduling, flight booking, airline crew paring, VLSI design, optimization for machine learning problems, etc \citep{TOC, CombApp1, CombApp2}. 

%No polynomial-time algorithm, which runs on a deterministic computing machine (any modern conventional computer say) and solves the problem exactly, is known for these problems as of now. Chances for such an algorithm to exist are very low because that will mean P = NP, which is very unlikely \citep{TOC}. Hence, in recent times, a lot of research has been devoted solving these problems heuristically on classical or emerging hardware. In many of these approaches, an Ising Hamiltonian is first formulated for a particular NP-hard problem \citep{NPCompleteIsingFormulation}, and then various techniques are used to reach the ground state of the Ising Hamiltonian, which is essentially the correct solution to the problem \citep{IsingComputingReview}. This emerging method of computing is known as Ising computing , and any hardware which implements it is called an Ising machine \citep{IsingComputingReview}.  

%A recent report by Mohseni \textit{et al.} \citep{IsingComputingReview} classifies Ising machines into three larger categories: classical annealers and their variations, quantum annealers and their variations like Quantum Approximation Optimization Algorithm (QAOA) \citep{QAOA_Farhi}, and dynamic-system solvers. In the classical annealing technique and its variations thereof, a fact from statistical physics is used: as temperature is lowered, the probability of a system to be in its lowest energy state grows \citep{IsingComputingReview}. Hence, in this technique, the temperature of the simulated system is lowered slowly to traverse over the energy landscape and reach the ground state. A major problem with this approach is the chance of getting stuck in a local minimum (wrong solution) instead of reaching the global minimum (correct solution). 

%On the other hand, quantum annealers (the second kind of Ising machines) try to reach the ground state from other states through quantum fluctuations \citep{IsingComputingReview}.  It has been shown that for some problem instances, where there are tall and narrow energy barriers separating the various minima in the energy landscape, quantum fluctuations can indeed enable the quantum annealer to find the ground state faster than its classical counterpart \citep{QuantumvsClassicalAnnealing4,QuantumvsClassicalAnnealing5,QuantumvsClassicalAnnealing6,QuantumvsClassicalAnnealing2,QuantumvsClassicalAnnealing3}. 


Noisy intermediate scale quantum (NISQ) hardware is another kind of unconventional hardware proposed for Ising computing. The current quantum era is deemed to be the NISQ era with all gate-based quantum hardware having 50-100 noisy qubits at the maximum \citep{NISQReview}. Since the originally proposed quantum algorithms like Shor's algorithm and Grover's search algorithm cannot be implemented on such a moderate number of noisy qubits (either very few noiseless qubits or millions of noisy qubits are needed for their implementation), various quantum-classical hybrid algorithms suitable for the NISQ era have evolved. The quantum approximate optimization algorithm (QAOA) is one such NISQ-era algorithm. Since the quantum circuit for QAOA has shallow depth and a classical optimization algorithm is used to optimize the quantum circuit parameters (more details on this in Section II of the paper), QAOA has been shown to be very suitable for NISQ-era quantum machines \citep{NISQReview, QAOA_Farhi,QAOA_Tutorial}. 

However, the quantum circuit involved in QAOA still needs very low temperature (in milliKelvins) for operation since quantum coherence needs to be maintained for the qubits during each forward computation step using this circuit (more details in Section II). Compared to that, PBOs are known to operate completely at room temperature since the governing physics of auto-oscillation and coupling do not involve quantum coherence. Hence, PBOs are more suitable compared to QAOA hardware for building practical machines for Ising computing (also known as Ising machines). Additionally, if we can show that the performance of PBOs is better than that of QAOA for different combinatorial optimization problems, then we show an added advantage of PBOs over QAOA.  

Motivated by this factor, here we have compared the performance of PBOs, simulated using Kuramoto model, with that of QAOA, simulated on the Qiskit simulator, for various instances of the popular NP-Hard combinatorial optimization problem: MaxCut (defined in Section II) \citep{IsingComputingReview,Spin_Akerman3,MIT2}. To cover instances of different levels of difficulty, we use four kinds of graphs: three of unweighted kinds (Mobius ladder, random cubic, Erd{\"o}s R{\'e}nyi) and one of weighted kind (complete graphs) (more details in Appendix A). We use graph size (number of nodes) from 4 to 20. We show that indeed PBOs exhibit much superior performance compared to QAOA for a wide range of graph instances. In fact, the success probability (probability to find the correct Max-Cut solution) of PBOs is 4-5 orders of magnitude above that of QAOA for the more difficult problem instances that we have considered here: random cubic, Erd{\"o}s R{\'e}nyi, and complete graphs with relatively high number of nodes: 18-20 (Fig. ~\ref{SuccessProb}). This shows that PBOs can indeed provide a more suitable platform compared to QAOA for practical Ising computing.

Since Kuramoto model is a generalized physics-agnostic model that can describe the dynamics of a wide range of oscillators (as mentioned above), our result will be applicable to a wide range of oscillators: metal-to-insulator-transition-based \citep{MIT1,MIT2}, spintronics-based \citep{Spin_Akerman1,Spin_Liu,Spin_Finocchio1,Spin_Finocchio2}, analog-and-digital-circuits based \citep{Circuit1,Circuit2,Circuit3,Circuit4,Circuit5}, etc. But the time to solution for different PBO-based Ising machines will be different; it will be inversely proportional to the natural frequency of operation of the particular oscillator devices in use. So in this paper, we have also reported our estimated time to solution (for the different graphs we considered here) for these different oscillator technologies. 





%To the best of our knowledge, this kind of comparison study between PBO and QAOA has not been reported before for such a wide range of Max-Cut problem instances   

%The quantum app QAOA has rapidly emerged as an alternative to quantum annealing.  While quantum annealing (typically implemented on annealers like D-Wave Systems) takes a long time to reach the ground state because of its adiabatic nature, QAOA (which is a Trotterized version of quantum annealing) can be implemented on shallow-depth quantum circuits (gate-based quantum computing model). Hence, QAOA is highly suitable for the current noisy-intermediate-scale-quantum (NISQ) era of quantum computing \citep{QAOA_Farhi,QAOA_Tutorial,NISQReview}. Having high degree of qubit-to-qubit connectivity is also an additional challenge in quantum annealing (D-Wave Systems), making QAOA quite popular. The performance of QAOA has hence been benchmarked against quantum annealing \citep{QuantumAnnealingvsQAOA1,QuantumAnnealingvsQAOA2,QuantumAnnealingvsQAOA3} and with classical annealing and other classical optimization algorithms \citep{QuantumAnnealingvsQAOA2,QAOAvsClassicalAnnealing1,QAOAvsClassicalAnnealing2} for different problems and problem instances. 


%However, the major disadvantage with both quantum annealing and QAOA is that they need very low temperature (in mK) for operation. Compared to that, dynamic-system solvers mostly (the third kind of Ising machines \citep{IsingComputingReview}) operate at room temperature, which is a major advantage for these systems. However, there are very few reports published thus far which have compared the performance of any dynamics-system-based Ising hardware with that of quantum Ising hardware. Hamerly \textit{et al} have recently compared the performance of coherent Ising machines, in the form of optical-oscillator-based dynamic systems, with that of quantum annealers for the same optimization problem instances \citep{CoherentIMvsQA}. They have shown experimentally that the optical-oscillator-based dynamic system (also known as coherent Ising machine, or CIM) exhibits superior performance compared to the D-Wave Systems quantum annealer; compared to D-Wave, CIM reportedly has orders of magnitude lower time to solution (TTS) and orders of magnitude higher probability of finding the ground state (success probability) for complex problem instances \citep{CoherentIMvsQA}. But CIM has major disadvantages like the use of a long fibre ring cavity and a power hungry field-programmable gate array (FPGA) to implement coupling \citep{MIT2}.  

%As an alternative to CIM, another kind of oscillator-based hardware has recently emerged as a dynamic system for Ising computing, which also operates at room temperature like CIM. Here, the coupled oscillators follow  the synchronization dynamics predicted by Kuramoto model \citep{IsingComputingReview,Kuramoto}, or that by the generalized Adler's equation (Gen-Adler)  \citep{Jaijeet1,Jaijeet2}. Upon the application of an external signal of frequency twice the average natural frequency of the oscillators (sub-harmonic injection locking, or SHIL), their phases settle down to either odd or even multiples of $\pi$, also known as phase binarization under  \citep{Jaijeet1,Jaijeet2}. Oscillators built from various kinds of conventional and emerging devices, including analog and digital electronic oscillator circuits \citep{Circuit1,Circuit2,Circuit3,Circuit4,Circuit5}, metal-to-insulator-transition-based oscillators \citep{MIT1,MIT2}, spintronic oscillators \citep{Spin_Jullie1,Spin_Jullie2,Spin_Akerman1,Spin_Akerman2,Spin_Akerman3,Spin_Liu,Spin_Hyunsoo,Spin_Ashwin,Spin_DB1,Spin_DB2},  etc., have been shown to follow the synchronization dynamics predicted by Kuramoto model or Gen-Adler equation. These oscillators have also been shown to exhibit phase binarization due to sub-harmonic injection locking (SHIL), and as a result Ising computing has been demonstrated experimentally on a network of such oscillators \citep{Circuit1,Circuit2,Circuit3,Circuit4,MIT2,Spin_Akerman1}. The emerging-nano-device-based oscillators, like that based on metal-to-insulator transition \citep{MIT1,MIT2} or spintronics \citep{Spin_Akerman1, Spin_Liu}, offer the advantage of high scalability and low energy consumption unlike CIM \citep{CoherentIMvsQA,MIT2}. Hence, these oscillator systems have been heavily pursued for Ising computing. But to the best of our knowledge, the performance of such phase-binarized-oscillator-based hardware has not been compared with that of quantum hardware thus far for the exact same problem instances (and also covering a wide range of problem instances), unlike in the case of CIM \citep{CoherentIMvsQA}. Even CIM's performance has only been compared with that of D-Wave-Systems-based quantum annealing, and not with QAOA (which has several advantages over quantum annealing, as mentioned above).  


%So, in this paper, we have compared the performance of a system of phase-binarized oscillators with that of QAOA for Ising computing over a wide range of problem instances, while restricting ourselves to one particular NP-Hard combinatorial optimization problem only: Max-Cut. We use four kinds of graphs (Mobius Ladder, random cubic, Erdos Renyi) to generate this wide range of problem instances. We use the Kuramoto model to model the phase-binarized oscillators; since Kuramoto model is a generalized physics-agnostic model that can describe the dynamics of a wide range of oscillators (as mentioned above), our study should be applicable to a wide range of oscillators: metal-to-insulator-transition-based \citep{MIT1,MIT2}, spintronics-based \citep{Spin_Akerman1,Spin_Liu}, analog-and-digital-circuits based \citep{Circuit1,Circuit2,Circuit3,Circuit4,Circuit5}, etc. 


%Among various quantum optimization algorithms, we have chosen QAOA because of its compatibility with gate-based quantum computing (like that implementable on various publicly available quantum hardware, e.g, IBMQ, Rigetti, etc). Since QAOA is a quantum-classical hybrid algorithm, it can be implemented on Noisy Intermediate Scale Quantum (NISQ) computers available currently; this has made QAOA a very popular quantum algorithm. 




%The classical annealing algorithm we have chosen is ..... QAOA is a Trotterized version of quantum annealing, which is based on the gate-based model of quantum computing. 


%To implement the classical annealing algorith, we have ..........


%To model the QAOA circuit and implement the corresponding algorithm, we have used the quantum circuit simulator Qiskit made available by IBM. To simulate the phase-binarized-oscillator hardware, we have used two different models. The first model is known as the Kuramoto model. It is hardware-agnostic and can be applied to a wide range of oscillators (chemical, mechanical, electronic, spintronic, etc.). It captures the auto-oscillation and synchronization of oscillators mathematically without taking into account the operating physics of the oscillators. The second model, formulated by Slavin and colleagues, is specific to spintronic nano-oscillators, which have been shown both theoretically and experimentally to exhibit auto-oscillation, synchronization, and phase binarization through sub-harmonic injection locking (SHIL). Spintronic nano-oscillators are particularly interesting among all possible oscillator technologies because of their very small size (in tens or hundreds of a nano-meter), high operating frequency (in GHz), electrical modulation of their natural frequencies and coupling coefficients, etc \citep{Akerman1,Akerman2}. The model proposed by Slavin and colleagues, referred to as Slavin's model, captures the coupled magnetization dynamics of spintronic oscillators by making approximations on the Landau Lifschitz Gilbert Slonczweski (LLGS) equations, as described in details in \citep{Slavin1,Slavin2}. \textbf{We may remove the Slavin's model from this paper because it is too hardware-specific if Kuramoto model gives better results (more details in Kuramoto sub-section of Section II).}


%In this paper, we compare the performance of QAOA and phase-binarized oscillators for the NP-Hard combinatorial optimization problem Max-Cut. To cover graph instances of a wide range of difficulty, we show solutions for four different kinds of graphs: Mobius Ladder, random cubic, Erdos Renyi, and weighted complete graphs. We compare the performances of QAOA and oscillators with each other and also a benchmark classical brute-force solver. We conclude that.....  



%Given that all Ising machines can be broadly classified into these three kinds, by choosing one algorithm each of one kind here, we have tried to cover the full spectrum of Ising machines through our paper. To the best of our knowledge, this is the first report of comparison of performance of three kinds of Ising machines (classical annealing, quantum, oscillator-based) for the same Max-Cut problem instances, which leads to a fair comparison. 



In Section II, we first formulate the Max-Cut problem and then provide details on the method we have followed to  model the system of PBOs (Kuramoto model) and implement QAOA (using Qiskit simulator provided by IBM and a classical optimization algorithm). In Section III, we compare the performance of QAOA and PBOs for the Max-Cut problem, solved using the aforementioned four kinds of graphs. In Section IV, we show the time-to-solution numbers (to solve the Max-Cut problem) for different device technologies considered for PBOs thus far. In Section V, we conclude the paper. 

\section{Problem Statement and Methods}

\begin{figure}[!t]
     \centering     \includegraphics[width=\columnwidth]{MobiusLadderInstances.jpg}
         \caption{(a) 6-node Mobius ladder and 8-node Mobius ladder graphs (unweighted) are shown without showing any partitioning. (b) Partitioning of the graphs in (a) such that the corresponding score for the Max-Cut problem (see description in Sec. IIA) is maximized is shown here. Nodes belonging to one partition are shown as dark blue circles, nodes belonging to the other partition are shown as yellow circles, and the edges which contribute to the final score that is to be maximised are shown as red dashed lines. The corresponding maximinized score (Max-Cut score) is written below each partitioned graph.}
    \label{MobiusLadder}
\end{figure}






\subsection{The Max-Cut Problem}

We define the Max-Cut problem formally as follows. Let there be a graph $G=(V,E)$ where $V$ is the set of vertices $\{1,2,3,...N\}$ and $E$ is the set of edges. Each member of E can be written as $(i,j)$, where $i \in V$ and $j \in V$ and $i < j$, with the mapping function C assigning a weight to each member of E. For an unweighted graph, $C: E \rightarrow \{0,1\} $, while for a weighted graph, $C: E \rightarrow R^{+}$. 

The problem statement of Max-Cut is that one has to select a sub-set $S$ of V ($S \subset G$) such that the following sum is maximized: $\Sigma_{e \in \delta_G(S)} C(e) $, where $\delta_G(S)=(i,j)$ with $i \in S$ and $j \in (V-S)$, or $i \in (V-S)$ and $j \in S$. 

Intuitively, when the graph is unweighted, the graph needs to be partitioned into two parts such that the number of edges which connect a node in one partition with a node in the other partition is maximised. This is known as the maximum cut of the graph. Such partitioning has been shown for a couple of Mobius Ladder graph instances (more details on the graph in Appendix A) in Fig. ~\ref{MobiusLadder}(b), with the same graphs without partitioning shown in Fig. ~\ref{MobiusLadder}(a).     


To cover a wide spectrum of problem instances and to check the efficiency of all the Ising computing algorithms we work on here on these instances, we use four kinds of graphs in this paper including Mobius ladder, as mentioned in Section I. More details on these four kinds of graphs are provided in Appendix A, along with images of some graph instances we have generated for this work (Fig. ~\ref{RandomCubic}, ~\ref{ErdosRenyi}).



\begin{figure*}[!t]
     \centering
\includegraphics[width=0.8\textwidth]{SchematicPBO.jpg}
     \caption{Schematic showing update of the phases of the oscillators $\phi_1$, $\phi_2$,....$\phi_N$ after every time step ($\Delta t$) in the phase-binarized oscillator (PBO) system, as per equation ~\ref{kuramoto_eqn} (Kuramoto model), as described in text. Total number of oscillators $N$ is same as the number of nodes for a given graph instance. The natural frequency values of all  oscillators  are considered the same here, making  $\omega _{i}-\omega^{*}$ term zero for any oscillator $i$ in equation ~\ref{kuramoto_eqn}. The coupling term between oscillator $i$ and oscillator $j$ corresponding to the edge connecting node $i$ with node $j$ of the graph is shown. For an unweighted graph, if there's an edge between node $i$ and node $j$, a constant coupling term is present between oscillator $i$ and oscillator $j$. For a weighted graph, the coupling between two oscillators $i$ and $j$ is proportional to the weight of an edge connecting node $i$ with node $j$. The SHIL term in equation ~\ref{kuramoto_eqn}, which leads to phase binarization, is also shown.}
     \label{PBO_Schematic}
\end{figure*}

\subsection{Modelling of Phase Binarized Oscillators (PBOs)}

As mentioned in Section I, Kuramoto model is a generalized, physics-agnostic model used to describe phase dynamics of a wide range of coupled oscillators. Here, we first briefly describe the model and then discuss how we have used it to model PBOs. We also provide the simulation parameters we used. 

According to the Kuramoto model, the net energy of a system of coupled oscillators, under sub-harmonic injection locking (SHIL), which leads to phase binarization, is given by the following equation:

\begin{equation}
E=-\Sigma_{i=1}^{n}\Sigma_{j=1,i < j}^{n}K_{ij}\cos(\phi_{i}-\phi_{j})-K_{s}\sum_{i}^{n}\cos (2\phi _{i})
\label{kuramoto_hamiltonian}
\end{equation}
, where $\phi_i$ is the phase of oscillator $i$, $K_s$ is the anisotropy constant (the term above with $K_s$ as coefficient leads to SHIL), and $K_{ij}$ is the coupling constant between oscillator $i$ and oscillator $j$ \citep{Spin_Liu}. When a given graph of $N$ nodes in the Max-Cut problem is mapped to a system of $N$ oscillators, the value of $K_{ij}$ is chosen to be proportional to the weight of the edge connecting nodes $i$ and $j$ in the corresponding graph. For the unweighted graphs that we consider here, when there's an edge between two nodes $i$ and $j$ of the graph, $K_{ij}=-K_{coupling}$ ($K_{coupling} >0$), and when there is no edge between nodes $i$ and $j$, $K_{ij} = 0$. For the weighted graphs we consider here, $K_{ij}=- w_{ij}K_{coupling}$, where $w_{ij}$ is the weight of the edge connecting nodes $i$ and $j$. The phases of the oscillators evolve over time ($t$) as follows:

\begin{equation}
\frac{d \phi_i}{dt}=\omega _{i}-\omega^{*} - \sum _{j,i\ne j}^{N}K_{ij}\sin (\phi _{i}-\phi_{j})-2K_{s}\sin (2\phi_{i})
\label{kuramoto_eqn}
\end{equation}

, where $\omega_i$ is the natural frequency of oscillator $i$ and $\omega^{*}$ is the average natural frequency of all the oscillators \citep{Spin_Liu}.

When the natural frequencies of the oscillators are equal (as in the cases considered in this paper), $\frac{d \phi_i}{dt} = -\frac{\partial E}{\partial \phi_i}$, and thus the phases evolve such that the net energy of the system $E$ (as per equation ~\ref{kuramoto_hamiltonian}) goes to a minimum. Once steady state is reached, oscillators with phases given by $2m\pi$ (for some $m\in \mathbb{Z}$) correspond to nodes in the graph belonging to one partition, and the oscillators with phases given by $(2m+1)\pi$ (for some $m\in \mathbb{Z}$) correspond to the nodes in the graph belonging to the other partition. From this partitioning, the corresponding cut value of the solution is obtained. 

%In Section III, we present the results we obtain this way for all the graphs we work on; we show the final output of the QAOA solver in each case, the evolution of the expectation value $\bra{\psi}H_P\ket{\psi}$ as a function of iteration number, and the measurement statistics (out of 1024 trials) using the forward-pass quantum circuit with updated  ($\beta_1$, $\gamma_1$, $\beta_2$, $\gamma_2$, ... $\beta_p$, $\gamma_p$) parameters. 





In our implementation of the Kuramoto model, we use the following set of parameter values: $\omega_i$ =0.1 radians/s for all oscillators labelled $i$ from 1 to $N$ (hence, $\omega^{*}$ = 0.1 radian/s),
$K_{coupling}=$ 1 radians/s, and $K_{s}=0.5$ radians/s. For each graph instance, we solve the coupled equations for oscillators (equation ~\ref{kuramoto_eqn}) numerically 100 times (100 trials) with different values of initial phases each time and thus obtain 100 solutions, some of which may coincide. The natural frequency values of all  oscillators  are considered the same here, making  $\omega _{i}-\omega^{*}$ term zero for any oscillator $i$. In each trial, we have used time-step $\Delta t=5 \times 10^{-5}$ s, and number of time-steps =$10^{6}$. Numerical computation of oscillator phases at time $t+\Delta t$ given phases at time $t$, as per equation ~\ref{kuramoto_eqn}, is shown in Fig. ~\ref{PBO_Schematic}.


For each trial, the final oscillator phase values are used to obtain the partitioning and cut value of the graph instance yielded for that trial, as described above. We use the results obtained in the 100 trials to report the cut value statistics in Section III. We also report there approximation ratio and success probability for each graph instance (the method for calculating approximation ratio 
and success probability is described in Section III).


%Among 100 trials, the solution with the highest cut value is used to calculate and report the approximation ratio of PBO. But to calculate success probability, 



%In the next section (Section III), we present the results we obtain this way through the Kuramoto model for all the graphs mentioned in Section IIB. For all the graph instances, we show the final output of the Kuramoto oscillator solver and compare it against that given by the QAOA solver and the classical benchmark solver. For proper comparison with QAOA, for the same graph instances where we show expectation value evolution and measurement statistics for QAOA, we also show the evolution of energy of oscillators (given by equation 4) vs iterations for the Kuramoto-oscillator case and the solution statistics (out of 100 trials) for the Kuramoto-oscillator case. 


\begin{figure*}[!t]
     \centering
    \includegraphics[width=\textwidth]{SchematicQAOA.jpg}
     \caption{Schematic of an implementation of Quantum Approximation Optimization Algorithm (QAOA), showing the quantum circuit for forward pass and the classical feedback/ optimization carried out to update the parameters of the quantum circuit: ($\beta_1$, $\gamma_1$, ... $\beta_p$, $\gamma_p$). In the quantum circuit, total number of qubits $N$ is same as the number of nodes for a given graph instance. The single-qubit gate H corresponds to the Hadamard gate. Unitary operator $e^{-i\beta H_M}$ can be decomposed into a single-qubit $R_X(2\beta)$ gate applied on each qubit. For an unweighted graph, unitary operator $e^{-i\gamma H_P}$ can be decomposed into two-qubit $R_{ZZ}(2\gamma)$ gates applied on every two qubits $i$ and $j$ if an edge connects node $i$ and node $j$, and no gate on qubits $i$ and $j$ if there's no edge between node $i$ and node $j$. For a weighted graph, $e^{-i\gamma H_P}$ can be decomposed into two-qubit $R_{ZZ}(2\gamma w_{i,j})$ gates applied on every two qubits $i$ and $j$, where $w_{i,j}$ is the weight of an edge connecting node $i$ and node $j$ of the graph.}
    \label{QAOA_Circuit}
\end{figure*}



\subsection{Numerical Implementation of Quantum Approximate Optimization Algorithm (QAOA)}


Here, we first briefly describe QAOA, in relation to solving MaxCut, and then discuss the numerical method we have used to implement QAOA. The MaxCut problem of any unweighted graph of $N$ nodes is solved through QAOA in the Hilbert space of $2^{N}$ dimensions, with each basis state corresponding to a possible partitioning of the graph. Thus solving this problem through a quantum circuit involves $N$ qubits. The given graph is converted to the corresponding problem Hamiltonian ($H_{P}$) in the same Hilbert space as follows \citep{QAOA_Farhi, QAOA_Tutorial}: 

\vspace{-2em}
\begin{multline}
H_P= \Sigma_{\textit{edge connecting node i and j}}(\frac{1}{2}w_{i,j}((I^{1} \\
\otimes I^{2} \otimes ... I^{i}\otimes...I^{j}\otimes...I^{N}) \\
-(I^{1} \otimes I^{2} \otimes .... \sigma_z^{i}\otimes ...\sigma_z^{j}\otimes... I^{N})))
\label{QAOA_Hp}
\end{multline}, 
where $\sigma_z$ corresponds to the Pauli spin matrix (in z) and $I$ corresponds to Identity matrix. As the first step of forward pass of every iteration, the initial state $\ket{s}$ is freshly prepared by applying Hadamard gate on all the $N$ qubits initialized to $\ket{0}$ each (Fig. ~\ref{QAOA_Circuit}) \citep{QAOA_Tutorial}. Then, as shown in Fig. ~\ref{QAOA_Circuit}, we let the state $\ket{s}$ evolve through multiple stages of unitary operators involving mixing Hamiltonian ($H_M$) and problem Hamiltonian ($H_P$) (defined above) and coefficients/ learning parameters ($\beta_1$, $\gamma_1$, $\beta_2$, $\gamma_2$, ... $\beta_p$, $\gamma_p$; if there are total $p$ stages), as follows \citep{QAOA_Farhi, QAOA_Tutorial}: 

\begin{multline}
    \ket{\psi}=e^{-i\beta_pH_M}e^{-i\gamma_pH_P}.....\\
    e^{-i\beta_2H_M}e^{-i\gamma_2H_P}e^{-i\beta_1H_M}e^{-i\gamma_1H_P} \ket{s}
    \label{QAOA_evol}
\end{multline}. This is the set of unitary operations corresponding to the forward-pass quantum circuit here. The mixing Hamiltonian $H_M$ is given by: 

\begin{equation}
H_M=\sum _{i=1}^{N}(I^{1}\otimes I^{2} \otimes I^{3}...I^{i-1} \otimes \sigma _{x}^{i}\otimes I^{i+1}....\otimes I^{n})
\label{QAOA_Hm}
\end{equation}. 

The gate decomposition of the unitary operators $e^{-i\gamma H_P}$ and $e^{-i\beta H_M}$ is shown in Fig.~\ref{QAOA_Circuit}. $e^{-i\beta H_M}$  can be decomposed into a set of $N$ single-qubit $R_X(2\beta)$ gates, with each gate applied on each one of the $N$ qubits. For an unweighted graph, $e^{-i\gamma H_P}$ can be decomposed into a set of two-qubit $R_{ZZ}(2\gamma)$ gates, with each such $R_{ZZ}$ gate applied on every two qubits $i$ and $j$ if an edge connects node $i$ and node $j$; no such gate is applied on qubits $i$ and $j$ if there's no edge between node $i$ and node $j$. For a weighted graph, $e^{-i\gamma H_P}$ can be decomposed into two-qubit $R_{ZZ}(2\gamma w_{i,j})$ gates applied on every two qubits $i$ and $j$, where $w_{i,j}$ is the weight of an edge connecting node $i$ and node $j$ of the graph.


As shown in Fig. ~\ref{QAOA_Circuit}, the expectation value of the problem Hamiltonian for this evolved state ($\bra{\psi}H_P\ket{\psi}$) is maximized over several iterations, where each update consists of a forward pass following equation 2 and an update  of ($\beta_1$, $\gamma_1$, $\beta_2$, $\gamma_2$, ... $\beta_p$, $\gamma_p$) following some classical optimization algorithm which makes use of the expectation value $\bra{\psi}H_P\ket{\psi}$ \citep{QAOA_Farhi,QAOA_Tutorial}.  If we consider the negative of the expectation value, then we are minimizing this value and reaching the ground state of the corresponding Hamiltonian, or the global minimum of the corresponding energy landscape \citep{IsingComputingReview}.

 

In our paper, we have used Python 3 with IBM's Qiskit quantum simulation package \citep{Qiskit} to implement QAOA. The forward computation, as given by equation ~\ref{QAOA_evol} and as shown in Fig. ~\ref{QAOA_Circuit}, is simulated using the Qiskit circuit simulator. We utilized the Aer Simulator in Qiskit with default parameters and without any noise model to simulate QAOA on an ideal quantum computer. Under the default settings, the simulator automatically selects the best method for the implementation of quantum circuit and its measurement. To keep the simulation time reasonably low and yet have optimized performance, we have used 10 stages only in our QAOA circuit of Fig. ~\ref{QAOA_Circuit} for all the graph instances we use ($p$=10). Nelder-Mead optimization algorithm \citep{NelderMead}, provided in the SciPy package, has been used to optimize the expectation value ($\bra{\psi}H_P\ket{\psi}$) through the classical feedback loop (Fig. ~\ref{QAOA_Circuit}). For this, the values of ($\beta_1$, $\gamma_1$, $\beta_2$, $\gamma_2$, ... $\beta_p$, $\gamma_p$) are updated over several iterations as long as the difference in expectation value between two consecutive iterations is greater than $10^{-4}$. We choose Nelder-Mead method over other optimization methods because it has been shown previously in several reports that Nelder-Mead method \citep{NelderMead} is a very effective optimization method for this kind of problems \citep{QAOA_ErdosRenyi, QuantumAnnealingvsQAOA1,QAOA_NelderMead1,QAOA_NelderMead2}. 


Before the first iteration, initial values of ($\beta_1$, $\beta_2$, ... $\beta_p$) are sampled uniformly at random from [0,$\pi$] because $e^{-i\beta H_M}$ is a periodic function of $\beta$ with a periodicity of $\pi$. Initial values of ($\gamma_1$, $\gamma_2$, ... $\gamma_p$) are sampled uniformly at random from [0,$2\pi$] because $e^{-i\gamma H_P}$ is a periodic function of $\gamma$ with a periodicity of $2\pi$ \citep{QAOA_Farhi,QAOA_ErdosRenyi,QAOA_CompleteGraph,QAOA_NelderMead1,QAOA_NelderMead2}. 

Once the difference between the expectation value ($\bra{\psi}H_P\ket{\psi}$) is less than $10^{-4}$, the iterative parameter update process ends and we obtain the final values of ($\beta_1$, $\gamma_1$, $\beta_2$, $\gamma_2$, ... $\beta_p$, $\gamma_p$). With these final circuit parameter values, the forward  pass and subsequent measurement (corresponding to $\sigma_z$ operator) are carried out 1024 times (sometimes, it's carried out a million times, as mentioned in Section III). For each measurement outcome, the corresponding graph partitioning is obtained the same way it has been carried out for PBOs; here in QAOA, qubit $i$ being $\ket{0}$ means the corresponding node $i$ belongs to the first partition, and qubit $i$ being $\ket{1}$ means the corresponding node $i$ belongs to the second partition. Based on the partitioning, the corresponding cut score is calculated like before. We present the cut value statistics obtained this way in Section III, where we also report approximation ratio and success probability of QAOA on every graph instance using this measurement statistics (the method for  calculating approximation ratio and success probability is described in Section III).


\begin{figure*}[!t]
     \centering
\includegraphics[width=0.99\textwidth]{ApproximationRatioFigure.jpg}
\caption{Approximation ratio as a function of number of nodes of the graph, when phase-binarized oscillators (PBO) are used to solve the Max-Cut problem (red plot), and when quantum approximate optimization algorithm (QAOA) is used for the same (green plot). Four kinds of graphs are used: (a) Mobius ladder, unweighted, (b) random cubic, unweighted, (c) Erd{\"o}s R{\'e}nyi, unweighted, (d) complete graph, weighted. The methodology for obtaining approximation ratio is described in the text (Section III).}
\label{ApproxRatio}
\end{figure*}



\begin{figure*}[!t]
     \centering
\includegraphics[width=0.99\textwidth]{Histogram16NodeFigure.jpg}
\caption{The ratio of number of trials that has yielded a given cut value to total number of trials (normalized counts) is plotted as a function of the cut value for (a) Mobius ladder, unweighted, with 16 nodes, (b) a particular 16-node instance of random cubic graph, unweighted, (c) a particular 16-node instance of Erd{\"o}s R{\'e}nyi graph, unweighted, (d) a particular 16-node instance of complete graph, weighted. For all the sub-figures, red bars correspond to trials using phase-binarized oscillators (PBOs), while green bars correspond to trials using quantum approximate optimization algorithm (QAOA).}
\label{Histogram_16Node}
\end{figure*}


\begin{figure*}[!t]
     \centering
\includegraphics[width=0.99\textwidth]{Histogram20NodeFigure.jpg}
\caption{The ratio of number of trials that has yielded a given cut value to total number of trials (normalized counts) is plotted as a function of the cut value for (a) Mobius ladder, unweighted, with 20 nodes, (b) a particular 20-node instance of random cubic graph, unweighted, (c) a particular 20-node instance of Erd{\"o}s R{\'e}nyi graph, unweighted, (d) a particular 20-node instance of complete graph, weighted. For all the sub-figures, red bars correspond to trials using phase-binarized oscillators (PBOs), while green bars correspond to trials using quantum approximate optimization algorithm (QAOA).}
\label{Histogram_20Node}
\end{figure*}







\begin{figure*}[!t]
     \centering
\includegraphics[width=0.99\textwidth]{SuccessProbabilityFigure.jpg}
\caption{Success probability (probability to find the correct Max-Cut solution) as a function of number of nodes of the graph, when phase-binarized oscillators (PBO) are used to solve the Max-Cut problem (red plot), and when quantum approximate optimization algorithm (QAOA) is used for the same (green plot). Four kinds of graphs are used: (a) Mobius Ladder, unweighted, (b) random cubic, unweighted, (c) Erd{\"o}s R{\'e}nyi, unweighted, (d) complete graph, weighted. The methodology for obtaining success probability is described in the text (Section III). Despite using 5,000,0000 trials, no trial has been successful for QAOA for Erd{\"o}s R{\'e}nyi
graph instances with 19 and 20 nodes.}
\label{SuccessProb}
\end{figure*}


\section{Performance Comparison: PBO vs QAOA}


\subsection{Approximation Ratio Comparison}


We first compare the approximation ratio obtained from our simulated PBO system and QAOA for the four kinds of graphs mentioned in Section II. To obtain the approximation ratio for PBOs, 100 trials are performed for every graph instance, with each trial being a distinct numerical simulation of equation ~\ref{kuramoto_eqn} with different phase initialization. From the stead-state phase values of the oscillators after every trial, the corresponding partitioning and cut value are obtained, as described in Section II. To obtain the approximation ratio, the highest cut value among the 100 trials is divided by the Max-Cut value obtained for the same graph instance through a brute-force classical Max-Cut solver we developed.


To obtain the approximation ratio for QAOA, for every graph instance, after the 
parameters of the quantum circuit ($\beta_1$, $\gamma_1$, $\beta_2$, $\gamma_2$, ... $\beta_p$, $\gamma_p$) reach their final values through an iterative process (as described in Section II), forward pass is carried out on the initial state $\ket{s}$  using the trained quantum circuit and measurement is carried out on the final state. This final process (forward pass using trained circuit and subsequent measurement) is repeated 1000 times. The partitioning and cut values for all measurement outcomes are computed, and the highest cut value obtained across the 1000 measurements is divided by the Max-Cut value obtained through a brute-force classical Max-Cut solver. The corresponding ratio obtained is reported as the approximation ratio for QAOA for that graph instance.

To calculate the approximation ratio both for PBOs and QAOA, we use the highest cut value obtained only (and ignore the other values) because by the definition of Max-Cut, no partitioning of the graph is possible with a higher cut value than the Max-Cut value. Hence, no final output basis state for QAOA is possible, or no list of final phases of oscillators is possible, for which the corresponding graph partition's cut value is higher than the Max-Cut value. Similar method is used to report the cut value for QAOA in the original proposal for QAOA by Farhi \textit{et al} \citep{QAOA_Farhi}.

Fig. ~\ref{ApproxRatio} shows the approximation ratio obtained for PBOs as a function of graph size (number of nodes) and that obtained for QAOA as a function of graph size for four kinds of graphs: Mobius ladder (Fig. ~\ref{ApproxRatio}(a)), random cubic (Fig. ~\ref{ApproxRatio}(b)), Erd{\"o}s R{\'e}nyi (Fig. ~\ref{ApproxRatio}(c)),  and complete graph (Fig. ~\ref{ApproxRatio}(d)). We observe that for Mobius Ladder graphs, both PBOs and QAOA yield accurate results for all graph sizes from 4 to 20 (approximation ratio 1 throughout). This is expected because Mobius Ladder forms a sub-class of easy instances within MaxCut and even has a polynomial-time solution \citep{MobiusLadder_PolySoln}. However, for the other kinds of graphs which are much more complicated and where polynomial-time solution is not known to exist, once the graph size (number of nodes) is above a certain threshold, approximation ratio for QAOA falls significantly below 1.

\subsection{Cut-Value Statistics and Success Probability Comparison}


But approximation ratio plots alone are not indicative of performance since we only consider the highest cut value amongst several trials to calculate the approximation ratio. While this method for calculating approximation ratio is justified (as discussed above), it's also important to look at the other trial outcomes and calculate success probability. The success probability of a solver is defined as the ratio of number of trials with the cut score equal to the correct Max-cut score obtained from the benchmark solver to the total number of trials. Without this analysis, one doesn't have an estimate of how many trials to perform to get the correct outcome, and in the extreme case, the number of trials needed will grow exponentially with graph size making the Ising computing technique as inefficient as the classical brute-force solver. 

As a result, in Fig. ~\ref{Histogram_16Node} and Fig. ~\ref{Histogram_20Node}, we compare the solution statistics, or cut-value statistics, for PBOs (across the 100 trials) with that for QAOA (across 1000 measurements) for graphs (of all four kinds) with 16 nodes and 20 nodes respectively. We choose these two particular graph sizes because Fig. ~\ref{ApproxRatio}(c) already shows lower approximation ratios for QAOA compared to PBOs for Erd{\"o}s R{\'e}nyi graphs of these graph sizes. From Fig. ~\ref{Histogram_16Node} and Fig. ~\ref{Histogram_20Node}, we observe the following:

\begin{enumerate}
    \item 16-node Mobius ladder graph: For PBOs, 95 out of 100 trials yield correct solution (success probability = 0.95), while for QAOA, only 633 out of 1024 trials yield correct results (success probability = 0.618) (Fig. ~\ref{Histogram_16Node}(a)). 

    \item 20-node Mobius ladder graph: 
    For PBOs, 85 out of 100 trials yield correct solution (success probability = 0.85), while for QAOA, 605 out of 1024 trials yield correct results (success probability = 0.59) (Fig. ~\ref{Histogram_20Node}(a)). 

    \item 16-node random cubic graph:  For PBOs, 55 out of 100 trials yield correct solution (success probability = 0.55), while for QAOA, only 85 out of 1024 trials yield correct results (success probability = 0.083) (Fig. ~\ref{Histogram_16Node}(b)).

    \item 20-node random cubic graph:  For PBOs, 60 out of 100 trials yield correct solution (success probability = 0.6), while for QAOA, only 7 out of 1024 trials yield correct results (success probability = 0.025) (Fig. ~\ref{Histogram_20Node}(b)).

    \item 16-node Erd{\"o}s R{\'e}nyi graph: For PBOs, 34 out of 100 trials yield correct solution (success probability = 0.34), while for QAOA, none of the 1024 trials yields correct results (Fig. ~\ref{Histogram_16Node}(c)). 

    \item 20-node Erd{\"o}s R{\'e}nyi graph: For PBOs, 30 out of 100 trials yield correct solution (success probability = 0.3), while for QAOA, none of the 1024 trials yields correct results (Fig. ~\ref{Histogram_20Node}(c)). 

    \item 16-node complete graph: Max-cut score for the weighted graph instance we have generated is 36.38. For PBOs, 14 out of 100 trials yield scores between 35 and 36. So, success probability here (defined as chance of having the solution roughly within + or - 0.5 or Max-Cut value) is 0.14. For QAOA, only 1 out of 1000 trials yields a score between 35 and 36 (success probability $\approx$ 0.001) (Fig. ~\ref{Histogram_16Node}(d)).

    \item 20-node complete graph: Max-cut score for the weighted graph instance we have generated is 55.29. For PBOs, 2 out of 100 trials yield scores between 55 and 56. So, success probability here is 0.02. For QAOA, none of the 1000 trials yields a score between 55 and 56 (Fig. ~\ref{Histogram_20Node}(d)).

    

    

\end{enumerate}


Thus, we see from the above examples that in most cases, success probability obtained for PBOs is significantly higher than that for QAOA. In fact, in  Fig. ~\ref{Histogram_16Node}(c) and Fig. ~\ref{Histogram_20Node}(c), (d), none of the 1000 trials for QAOA has led to the correct score, making it impossible to obtain QAOA's success probability for these graph instances with 1000 trials. So, in Fig. ~\ref{SuccessProb}, where we compare PBO's and QAOA's success probability as a function of graph size (number of nodes), we have used 1,000,000-5,000,000 trials for QAOA when working on the specific graph instances for which no successful result has been obtained with 1000 trials. 

Fig. ~\ref{SuccessProb}(a) shows that for Mobius Ladder graphs, the easy case among various Max-Cut graph instances, success probability using QAOA is of the same order of that using PBOs. But solving Max-Cut for the other types of graphs we use here is much more difficult than solving Max-Cut for Mobius Ladder \citep{MobiusLadder_PolySoln,QAOA_ErdosRenyi,QAOA_CompleteGraph}. Fig. ~\ref{SuccessProb} (c), (d) show that for Erd{\"o}s R{\'e}nyi and complete graphs, success probability of QAOA drops exponentially with increase in the number of nodes of the graph, while success probability of PBO only drops linearly with increase in the number of nodes. In fact, for random cubic, Erd{\"o}s R{\'e}nyi, and complete graphs of 18-20 nodes, success probability of QAOA is 4-5 orders lower compared to that of PBOs (Fig. ~\ref{SuccessProb} (b), (c), (d)). This proves our claim that for various complex instances of the Max-Cut problem, PBOs perform much better than QAOA in terms of success probability. 
It is to be noted here that despite using 5,000,0000 trials, we still have not obtained any successful trial for QAOA for Erd{\"o}s R{\'e}nyi
graph instances with 19 and 20 nodes. Hence, in Fig. ~\ref{SuccessProb}(c), success probability of QAOA is not plotted for these two graph instances. 



\section{Time to Solution for Different Oscillator Technologies}



Having shown that PBOs offer much higher success probability than QAOA for various Max-Cut problem instances, we now determine the time to solution (TTS) for these instances. In our modelling of PBOs (as described in Section II), the equation for the Kuramoto model (equation ~\ref{kuramoto_eqn}) is solved numerically using a fixed time step. Hence, TTS can be represented by the number of time steps. So, in Fig. ~\ref{TTS1}, we have plotted TTS (in terms of number of time steps) as a function of graph size/ number of nodes for the four kinds of graphs we have explored in this paper. The fact that we have conducted 100 trials for each problem instance is also accounted for in Fig. ~\ref{TTS1}. 

We observe that TTS doesn't monotonically increase with graph size for any kind of graph we have considered here. Though this trend is different from that reported for TTS in a few previous reports on PBOs \citep{Circuit3,MIT2}, it is not unexpected because the graph sizes considered here are relatively low compared to the other reports. Also, the other reports \citep{Circuit3,MIT2} follow an annealing scheme for PBO (the magnitude of the SHIL term in equation ~\ref{kuramoto_eqn}, $K_s$, gradually increases over time), while we follow a simpler scheme where $K_s$ is always constant and always high enough to cause SHIL. This may be another reason for the aforementioned discrepancy. Based on Fig. ~\ref{TTS1}, we report in Table ~\ref{TTS2} the average TTS (in terms of number of time steps) across all the graph instances from 4 nodes to 20 nodes for the four kinds of graphs we have considered here.

The actual TTS (in seconds), however, is the product of TTS (in terms of number of time steps) (as shown in Fig.~\ref{TTS1}) and the value of each time step. In our PBO model, each time step = $5\times 10^{-5}$ s since we have considered the average natural frequency of the oscillators $\omega^{*}$ to be 0.1 rad/s (or 0.016 Hz) and the coupling constant $K_{coupling}$ to be 1 radian/s (or 0.16 Hz) (as mentioned in Section II). Such value of coupling constant leads to synchronization range between two oscillators to be in the order of Hz only \citep{Spin_DB1,Vassileva}. However, when the PBO system is implemented physically through different oscillator technologies, the natural frequencies of the oscillators and synchronization ranges can vary from kHz (for electronic-circuit-based oscillators \citep{Circuit5,Circuit6}) to GHz (for spintronic oscillators \citep{Spin_Akerman1,Spin_DB1}). $\omega^{*}$ and $K_{coupling}$ in equation ~\ref{kuramoto_eqn} need to be scaled up accordingly, and the time step needs to be scaled down by the same proportionality factor. The actual TTS (in seconds) will go down accordingly; the higher the natural frequency and synchronization range of the oscillators, the faster do PBOs converge to the correct solution. Based on this analysis, in Table ~\ref{TTS3}, we report the natural frequency and synchronization range of various kinds of oscillator technologies (based on available simulation-based and experimental reports, which we have cited in Table ~\ref{TTS3}). Next to that, we provide the time step value (for equation ~\ref{kuramoto_eqn}) we have accordingly considered for these technologies.

Further in Table ~\ref{TTS3}, based on these time step values reported for different oscillator technologies and the average number of time steps (across all the graph instances from 4 nodes to 20 nodes) calculated from Fig. ~\ref{TTS1} (as listed in Table ~\ref{TTS2}), we have reported the average time (in seconds) that will be needed to solve Max-Cut problem for each kind of graph we have considered here for different oscillator technologies. We observe that spintronic oscillators offer lowest TTS due to their very high operating frequencies. But experimental demonstration of Ising computing in spintronic oscillators is only at a preliminary stage currently due to the emerging nature of this kind of device technology \citep{Spin_Akerman1}. On the other hand, electronic-circuit-based oscillators, though operating at a much lower frequency range, are very easy to fabricate, and large Ising networks of electronic oscillators have already been demonstrated experimentally \citep{Circuit1,Circuit2,Circuit3}.     



\begin{figure}[!t]
     \centering
\includegraphics[width=0.99\columnwidth]{TTS_timesteps_.jpg}
     \caption{Time to solution (TTS), in terms of number of time steps, as a function of the number of nodes of the graph, when phase binarized oscillators (PBO) are used to solve the Max-Cut problem. Results for all the four kinds of graphs considered in this paper are shown: Mobius ladder (green), random cubic (blue), Erd{\"o}s R{\'e}nyi (black), and complete graph (red).}
    \label{TTS1}
\end{figure}


\begin{table*}[!t]
\begin{tabular}{|c|c|} 
\hline
\textbf{Kind of Graph} & \textbf{ Average TTS (number of time steps)} \\
\textbf{} & \textbf{, across graph instances from 4 to 20 nodes} \\
\hline
Mobius ladder, unweighted & \textbf{$1.892 \times 10^{7}$} \\
Random cubic, unweighted  &  \textbf{$2.288 \times 10^{7}$} \\ 
Erd{\"o}s R{\'e}nyi, unweighted  &  \textbf{$1.639 \times 10^{7}$} \\
Complete graph, weighted & \textbf{$1.382 \times 10^{7}$} \\
\hline

\end{tabular}
\caption{Average TTS (in terms of number of time steps) across all the graph instances from 4 nodes to 20 nodes for the four kinds of graphs we have considered here}
\label{TTS2}
\end{table*}

\begin{table*}[t!]
\begin{tabular}{|c|c|c||c|c|c|c|c|} 
 \hline
\textbf{Oscillator} & \textbf{Natural} & \textbf{Synch.} &\textbf{Time Step} &
\textbf{Avg. TTS} & \textbf{Avg. TTS} & \textbf{Avg. TTS} & \textbf{Avg. TTS}
\\
\textbf{Technology} & \textbf{Frequency} & \textbf{Range} &\textbf{(s)} & \textbf{(s): Mobius} & \textbf{(s): random} & \textbf{(s): Erd{\"o}s} & \textbf{(s): complete}\\
\textbf{} & \textbf{} & \textbf{} &\textbf{} & \textbf{ladder} & \textbf{cubic} & \textbf{R{\'e}nyi} & \textbf{graph} \\
\hline
Electronic circuit: &  &  &  &  &  &  & \\
Wien Bridge oscillator \citep{Circuit5}   & 260 Hz \citep{Circuit5} & 95  Hz \citep{Circuit5} & $3.1\times 10^{-9}$   & $5.86 \times 10^{-2}$     & $7.07 \times 10^{-2}$   & $5.05 \times 10^{-2}$   & $4.28 \times 10^{-2}$ \\
\hline
Electronic Circuit: 
&  & not reported  &  &  &  &  & \\
LC oscillator \citep{Circuit3} & 50 kHz \citep{Circuit3} & in \citep{Circuit3} & $1.6\times 10^{-11}$  & $3.02 \times 10^{-4}$     & $3.65 \times 10^{-4}$   & $2.61 \times 10^{-4}$   & $2.21 \times 10^{-4}$ \\
\hline
Metal-to-insulator 
&  &  &  &  &  &  & \\
transition oscillator \citep{MIT1} & 30 kHz \citep{MIT1} & 20 kHz \citep{MIT1} & $2.65\times 10^{-11}$  & $5.01 \times 10^{-4}$     & $6.04 \times 10^{-4}$   & $4.32 \times 10^{-4}$   & $3.66 \times 10^{-4}$ \\
\hline 
Spintronic oscillator \citep{Spin_DB1} & 6.7 GHz \citep{Spin_DB1} & 0.2 GHz \citep{Spin_DB1} & $1.19\times 10^{-16}$  & $2.25 \times 10^{-9}$     & $2.71 \times 10^{-9}$   & $1.94 \times 10^{-9}$   & $1.64 \times 10^{-9}$ \\
\hline
\end{tabular}
\caption{Average TTS (in seconds) across all the graph instances from 4 nodes to 20 nodes for the four kinds of graphs we have considered here, for different kinds of phase binarized oscillator (PBO) technologies. For each PBO technology, the corresponding natural frequency, synchronization range, and time step (in seconds) needed to solve equation ~\ref{kuramoto_eqn} are also listed.}
\label{TTS3}
\end{table*}




\section{Conclusion}

Thus, in this paper, we have shown that phase binarized oscillators (PBOs) offer superior performance compared to quantum approximation optimization algorithm (QAOA) for a wide range of graph instances. In fact, the success probability for PBOs is several orders of magnitude higher compared to QAOA for more complex graph instances. With PBOs and QAOA being two popular approaches for carrying out Ising computing currently, our finding here is particularly important. In this paper, we have also estimated the time to solution for different oscillator technologies (electronic circuit based, metal-to-insulator transition based, spintronic) for the graph instances we have considered.


It is to be noted that another dynamic system approach exists for Ising computing apart from PBOs, which is known as coherent Ising machine (CIM) \citep{IsingComputingReview}. Recently, Hamerly \textit{et al} have compared the performance of CIM with that of quantum annealers for the same optimization problem instances (QAOA is an approximate gate-based variation of quantum annealing) \citep{CoherentIMvsQA}. They have shown experimentally that CIM exhibits superior performance compared to the D-Wave Systems quantum annealer; compared to D-Wave, CIM reportedly has orders of magnitude lower time to solution (TTS) and orders of magnitude higher success probability for complex problem instances \citep{CoherentIMvsQA}. Hence, our result here (success probability of PBO orders of magnitude higher than that of QAOA) is consistent with that reported by Hamerly \textit{et al} \citep{CoherentIMvsQA}, with PBO being similar to CIM (dynamic-system-based approach for Ising computing \citep{IsingComputingReview}) and QAOA being similar to quantum annealing (quantum approach for Ising computing \citep{IsingComputingReview}). Additionally, CIM has major disadvantages like the use of a long fibre ring cavity and a power hungry field-programmable gate array (FPGA) to implement coupling, which is not the case with PBO \citep{MIT2}. Hence, our result here that PBOs, just like CIM, exhibit much higher success probability compared to quantum approaches makes PBOs very attractive for Ising computing.

\section*{Acknowledgment}
We thank Fractal Analytics for providing financial support to this project and also thank Aniket Sadashiva, Siddhartha Santra, Himadri Shekhar Dhar, and Prateek Jain for their technical insights regarding this project.



\section*{Appendix A: Descriptions of the Types of Graphs Considered Here}



\begin{enumerate}
    \item \textbf{Mobius ladder graph:} 
     A Mobius Ladder graph (shown in Fig. ~\ref{MobiusLadder}(a)) of graph size $N$ ($N$ is a positive even integer) is an undirected unweighted 3-regular (also known as cubic) graph with $N$ nodes for which every node $i$ has edges of weight 1 with only three other nodes \citep{Ref_MobiusLadder}:
$$
\begin{cases}
i-1, i+1, (\frac{N}{2})+i, \hspace{1em} \text{if} \hspace{1em}  2 \leq i \leq \frac{N}{2} \\
i-1, i+1, i-(\frac{N}{2}), \hspace{1em} \text{if} \hspace{1em}  (\frac{N}{2})+1 \leq i < N \\
n, i+1, (\frac{N}{2})+i \hspace{1em} \text{if} \hspace{1em} i = 1\\
i-1, 1, i-(\frac{N}{2}) \hspace{1em} \text{if} \hspace{1em} i = N  \\
\end{cases}
$$

With all other nodes, edges that node $i$ shares have 0 weight, and hence those edges are not drawn in  Fig. ~\ref{MobiusLadder}, which shows Mobius ladder graphs of graph size (number of nodes) 6 and 8. In our paper, we work on Mobius ladder graphs with graph size $N$ (even integer) from 4 to 20.



\item \textbf{Random cubic graph:}  Any random cubic graph we use here is an undirected unweighted 3-regular graph where each node is randomly connected to three other nodes \citep{Ref_RandomCubic}. In our paper, for each even value of graph size $N$ from 4 to 20, we generate a random cubic graph using the NetworkX Python package \citep{SciPyProceedings_11, Ref_RandomCubic,gen_random_reg}. The specific 4-node, 8-node, 12-node, 16-node and 20-node random cubic graph instances we generate for this work are shown in Fig. ~\ref{RandomCubic}.


\item \textbf{Erd{\"o}s R{\'e}nyi graph:} Any Erd{\"o}s R{\'e}nyi graph we use here is given by $G(N,p)$ and is an undirected unweighted graph of $N$ nodes, where the probability of any edge between two nodes to exist (also known as edge probability) is independent of any other edge to exist and is given by $p$ ($0 \leq p \leq 1$) \citep{Ref_ErdosRenyi}. Thus, the expected number of nodes for this graph is $p$${N}\choose{k}$. In our paper, for each value of $N$ from 4 to 20, we generate an Erd{\"o}s R{\'e}nyi graph given by $G(N,0.5)$.  We choose edge probability $p$= 0.5 because it corresponds to Erd{\"o}s R{\'e}nyi graphs of the highest difficulty level in terms of solving the Max-Cut problem. Graphs with both very low level of connectivity (low value of $p$) and very high level of connectivity (high value of $p$) have been found to be much easier Max-Cut problem instances compared to that \citep{QAOA_ErdosRenyi}. The specific 4-node, 8-node, 12-node, 16-node and 20-node Erd{\"o}s R{\'e}nyi graph instances we generate for this work are shown in Fig. ~\ref{ErdosRenyi}. 


\item \textbf{Complete graph:} Any complete graph we use here is an undirected but weighted graph, where every node is connected to every other node by an edge. A weight is assigned randomly to every edge, independent of other edges, with each weight sampled uniformly at random from [0,1] \citep{Ref_CompleteGraph,QAOA_CompleteGraph}.


\end{enumerate}

\begin{figure}[!t]
     \centering
\includegraphics[width=0.8\columnwidth]{RandomCubicInstances.jpg}
     \caption{Random cubic (3-regular) graph instances for different graph sizes (number of nodes: 4, 8, 12, 16, 20) that have been generated and used for this study.}
    \label{RandomCubic}
\end{figure}


\begin{figure}[!t]
     \centering
\includegraphics[width=0.85\columnwidth]{ErdosRenyiInstances.jpg}
     \caption{Erd{\"o}s R{\'e}nyi graph instances for different graph sizes (number of nodes: 4, 8, 12, 16, 20) that have been generated and used for this study. Probability for edge connectivity is taken to be 0.5 here to cover the most difficult Erd{\"o}s R{\'e}nyi graph instances for the Max-Cut problem}
    \label{ErdosRenyi}
\end{figure}




%Each of these times, the corresponding state that is obtained is one of the basis states and corresponds to a particular partitioning of the graph. So for each such obtained state, we calculate the corresponding cut value. The highest cut value obtained this way, across 1024 trials, is used here as the final output of the QAOA max-cut solver to report the approximation ratio (following the method proposed by Farhi et al \citep{QAOA_Farhi}).



%But we use all the cut values obtained across all the trials to calculate the success probability of QAOA. 

%In the next section (Section III), we present the results we obtain this way for all the graphs mentioned in Section IIB. We show the final output of the QAOA solver for each graph instance and compare it against that yielded by the benchmark classical brute-force solver for the same graph instance. For many graph instances, we also show the evolution of the expectation value $\bra{\psi}H_P\ket{\psi}$ as a function of iteration number, and the measurement statistics (out of 1024 trials) using the forward-pass quantum circuit with updated  ($\beta_1$, $\gamma_1$, $\beta_2$, $\gamma_2$, ... $\beta_p$, $\gamma_p$) parameters. 










%corresponding evolved state is supposed to be the eigen state of the problem Hamiltonian $H_P$ with the lowest eigen value (the ground state); hence, it corresponds to the partitioning of the graph which results in the highest Max-Cut value (i.e. the number of edges for which one node belongs to one side of the cut and the other node belongs to the other side of the cut is maximum for this partitioning). 





%\subsection{\label{sec:level2}Second-level heading: Formatting}

%This file may be formatted in either the \texttt{preprint} or
%\texttt{reprint} style. \texttt{reprint} format mimics final journal output. 
%Either format may be used for submission purposes. \texttt{letter} sized paper should
%be used when submitting to APS journals.

%\subsubsection{Wide text (A level-3 head)}
%The \texttt{widetext} environment will make the text the width of the
%full page, as on page~\pageref{eq:wideeq}. (Note the use the
%\verb+\pageref{#1}+ command to refer to the page number.) 
%\paragraph{Note (Fourth-level head is run in)}
%The width-changing commands only take effect in two-column formatting. 

\newpage

%\bibliography{OriginalManuscript_PRApp}% Produces the bibliography via BibTeX.

\input{OriginalManuscript_PRApp.bbl}

\end{document}
%
% ****** End of file apssamp.tex ******
